\begin{lemma}
Let $H$ be a hypergraph and let $v_0,v_1,v_2$ cover the
vertices of an edge $f$. Let $S$ be a set of three edge indices, each
meeting $f$, and suppose distinct members of $S$ are nonadjacent in
the line graph. Then there is a map $q:\{0,1,2\}\to S$ whose image is
$S$, with $v_i\in H(q(i))$. Moreover, any member of $S$ containing
$v_i$ equals $q(i)$.
\end{lemma}
\begin{proof}
Choose for each member $e$ of $S$ a vertex of $H(e)\cap H(f)$ and an
index $r(e)$ representing it. Distinct members cannot have the same
index: they would share $v_{r(e)}$ and hence be adjacent by
\noderef{LineGraphOfHypergraph}. Thus $r$ is an injection between two
three-element sets, and is a bijection. Take its inverse $q$.
Its image is $S$ and $q(i)$ contains $v_i$. Any other member containing
$v_i$ would intersect $q(i)$, contradicting nonadjacency. This proves
the final uniqueness assertion.
\end{proof}
